\chapter{Практика}


\section{Настройка интерфейсов роутера}

\begin{lstlisting}[language=C, caption={Пример настройки интерфейсов роутера}]
#R1
configure

interface gigabitethernet 1/0/1
ip address 192.168.10.1/30
ip firewall disable
exit

interface gigabitethernet 1/0/2
ip address 10.10.10.1/24
ip firewall disable
exit

1interface gigabitethernet 1/0/3
ip address 172.16.10.1/24
ip firewall disable
exit

do commit
do confirm
 
#R2
configure

interface gigabitethernet 1/0/1
ip address 192.168.10.2/30
ip firewall disable
exit

interface gigabitethernet 1/0/2
ip address 192.168.100.1/30
ip firewall disable
exit

do commit
do confirm

#R-DHCP
configure

interface gigabitethernet 1/0/1
ip address 192.168.100.2/30
ip firewall disable
exit

do commit
do confirm
\end{lstlisting}

\section{Настройка DHCP-сервера на роутере}

\begin{lstlisting}[language=C, caption={Пример настройки DHCP-сервера}]
# первый пул
R-DHCP(config)# ip dhcp-server
R-DHCP(config)# ip dhcp-server pool DHCP-10
R-DHCP(config-dhcp-server)# network 10.10.10.0/24
R-DHCP(config-dhcp-server)# address-range 10.10.10.100-
10.10.10.200
R-DHCP(config-dhcp-server)# default-router 10.10.10.1
R-DHCP(config-dhcp-server)# default-lease-time 0:0:30
R-DHCP(config-dhcp-server)# exit

# второй пул
R-DHCP(config)# ip dhcp-server pool DHCP-172
R-DHCP(config-dhcp-server)# network 172.16.10.0/24
R-DHCP(config-dhcp-server)# address-range 172.16.10.100-
172.16.10.200
R-DHCP(config-dhcp-server)# default-router 172.16.10.1
R-DHCP(config-dhcp-server)# default-lease-time 0:0:30
R-DHCP(config-dhcp-server)# exit
\end{lstlisting}

\section{Настройка ip helper}

ip helper необходим потому, что dhcp-client отправляет пакет dhcp-discover, который является широковещательным. Поскольку
сервер находится в другой сети, пакет до него не дойдет (роутер не пропускает широковещательный трафик). Роутер должен принять 
dhcp-discover, преобразовать его в unicast пакет и направить к серверу

\begin{lstlisting}[language=C, caption={Пример настройки ip helper}]
# включаем dhcp-relay
R1(config)# ip dhcp-relay
R1(config)# interface gigabitethernet 1/0/2
#задаем адрес сервера
R1(config-if-gi)# ip helper-address 192.168.100.2
R1(config-if-gi)# exit

R1(config)# interface gigabitethernet 1/0/3
R1(config-if-gi)# ip helper-address 192.168.100.2
R1(config-if-gi)# exit

R1(config)# do commit
R1(config)# do confirm
\end{lstlisting}


\section{Настройка статического маршрута}

Чтобы роутер мог отправить dhcp-discover в другую сеть нужно настроить статический маршрут R1->R-DHCP через R2.

\begin{lstlisting}[language=C, caption={Пример настройки статического маршрута}]
R1(config)# ip route 0.0.0.0/0 192.168.10.2
R1(config)# do commit
R1(config)# do confirm

R-DHCP(config)# ip route 0.0.0.0/0 192.168.100.1
R-DHCP(config)# do commit
R-DHCP(config)# do confirm

R2(config)# ip route 10.10.10.0/24 192.168.10.1
R2(config)# ip route 172.16.10.0/24 192.168.10.1
R2(config)# do commit
R2(config)# do confirm

\end{lstlisting}


\section{Проверка сетевой связности}

Сделаем ping ПК1-ПК2  

\begin{figure}[H] 
    \centering
    \includegraphics[width=0.75\textwidth]{dhcp.png}
    \caption{DHCP-клиент}
\end{figure}

Выполнили настройку клиента с помощью dhcp и проверили сетевую связность. Пинг проходит успешно.










\endinput